% =====================================================================
\section{IoT: ESP-NOW, MQTT и Bluetooth LE}
% =====================================================================
\lessonintro
  {выбрать подходящий беспроводной протокол и обмениваться данными без блокировки устройства}
  {урок 8 (Wi-Fi), урок 7 (задачи и очереди)}
  {плата становится беспроводным датчиком, частью умного дома или устройством для настройки с телефона}
  {вторая ESP32-S3; для MQTT --- брокер; для BLE --- телефон}

\subsection{Три задачи --- три протокола}

ESP-NOW нужен для коротких сообщений между ESP без роутера: пульт, кнопка, распределённый датчик.
MQTT нужен, когда независимые устройства должны общаться через брокер и Home Assistant. BLE удобен для
настройки телефоном рядом с устройством и экономит энергию. Это не взаимозаменяемые технологии:
выбирайте протокол по пути данных, а не по модному названию.

\begin{center}\small
\begin{tabularx}{\textwidth}{@{}l X X@{}}
\toprule
Протокол & Хорош для & Главное ограничение \\
\midrule
ESP-NOW & коротких команд, малой задержки, без роутера & обе платы на одном Wi-Fi-канале \\
MQTT & умного дома, журналов, нескольких клиентов & нужен Wi-Fi и брокер \\
BLE GATT & настройки и данные рядом со смартфоном & небольшие пакеты, не «обычный Bluetooth Serial» \\
\bottomrule
\end{tabularx}
\end{center}

\subsection{ESP-NOW: пакет по MAC-адресу}

Отправитель добавляет приёмник как peer и посылает структуру фиксированного размера. В реальном
проекте добавьте версию протокола, счётчик \code{sequence} и контроль диапазонов на приёмнике.
Коллбэк «доставлено» означает подтверждение на уровне радио, но не заменяет подтверждение, что
приложение действительно обработало команду.

\codecap{Отправка структуры раз в секунду}
\codeinput{l12_espnow}

\begin{caution}
ESP-NOW и обычный Wi-Fi используют один радиоканал. Если плата подключилась к роутеру на канале 6,
вторая плата с ESP-NOW должна оказаться на том же канале. MAC-адрес печатается через
\code{WiFi.macAddress()}; не копируйте адрес из примера. В ядре 3.3 изменился первый
аргумент колбэка отправки, поэтому пример выбирает нужный вариант через \code{ESP\_ARDUINO\_VERSION}.
\end{caution}

\subsection{MQTT: темы вместо прямых соединений}

В MQTT устройство подключается к брокеру и публикует сообщения в \emph{topics}, например
\code{home/esp-01/state}. Другой клиент подписывается на эту тему; команда приходит в
\code{home/esp-01/set}. Для состояния удобно оставить retained-сообщение, а для команд --- нет.
Используйте компактный JSON, например \code{\{"on":true,"level":75\}}, и проверяйте каждое поле до
применения. Для Arduino обычно устанавливают библиотеку PubSubClient через Library Manager.

Подключение к брокеру не помещают в бесконечный цикл: выполняйте одну попытку по таймеру, увеличивайте
интервал после ошибок и давайте работать кнопке, экрану и \code{client.loop()}. Логин/пароль и TLS
не делают публичный брокер бесконечным ресурсом: соблюдайте его ограничения и не публикуйте данные
десятки раз в секунду без причины.

\begin{figure}[H]
\centering
\begin{tikzpicture}[font=\sffamily\scriptsize]
  \node[blk,fill=blkrf,minimum width=30mm,minimum height=12mm] (b) at (0,0) {MQTT-брокер\\ \scriptsize Mosquitto, Home Assistant};
  \node[blk,fill=blkcpu,minimum width=28mm,minimum height=12mm] (e) at (-7.3,0) {ESP32-S3\\ \scriptsize клиент \code{esp-01}};
  \node[blk,fill=blkper,minimum width=28mm,minimum height=10mm] (p) at (6,1.1) {Телефон};
  \node[blk,fill=blkper,minimum width=28mm,minimum height=10mm] (h) at (6,-1.1) {Панель умного дома};
  \draw[arr] ($(e.east)+(0,0.25)$) -- node[above]{\code{home/esp-01/state}} ($(b.west)+(0,0.25)$);
  \draw[arr] ($(b.west)+(0,-0.25)$) -- node[below]{\code{home/esp-01/set}} ($(e.east)+(0,-0.25)$);
  \draw[darr] (b.east) -- node[above,sloped,pos=0.45]{state / set} (p.west);
  \draw[darr] (b.east) -- node[below,sloped,pos=0.45]{state / set} (h.west);
\end{tikzpicture}
\caption{Клиенты MQTT не соединяются друг с другом напрямую: они публикуют сообщения в темы брокера
и подписываются на нужные темы}
\end{figure}

\codecap{MQTT-клиент: неблокирующее переподключение, подписка и публикация состояния}
\codeinput{l12_mqtt}


\subsection{BLE: сервисы и характеристики}

BLE-сервер публикует сервис (например, «настройки устройства»), а внутри --- характеристики: яркость,
состояние, имя. У каждой есть UUID и права read/write/notify. Стандартный минимальный MTU даёт всего
20 байтов полезных данных в одном уведомлении, поэтому длинный JSON нужно фрагментировать или
согласовать увеличенный MTU. Для нового проекта предпочтительна библиотека NimBLE: она экономнее
старого BLE-стека. В примере используется встроенная в Arduino-ESP32 библиотека \code{BLEDevice}:
ей не нужна отдельная установка.

\begin{figure}[H]
\centering
\begin{tikzpicture}[font=\sffamily\scriptsize,ch/.style={blk,fill=blkmem,minimum width=30mm,minimum height=11mm}]
  \node[blk,fill=blkcpu,minimum width=40mm] (srv) at (0,0) {GATT-сервер ESP32-S3};
  \node[blk,fill=blkrf,minimum width=40mm] (svc) at (0,-1.4) {Сервис «Настройки»\\ \scriptsize UUID 4fafc201-\ldots};
  \node[ch] (c1) at (-4,-3.1) {Яркость\\ \scriptsize read, write};
  \node[ch] (c2) at (0,-3.1) {Состояние\\ \scriptsize read, notify};
  \node[ch] (c3) at (4,-3.1) {Имя\\ \scriptsize read, write};
  \draw[arr] (srv) -- (svc);
  \foreach \c in {c1,c2,c3} \draw[arr] (svc.south) -- (\c.north);
\end{tikzpicture}
\caption{Устройство BLE описывает себя деревом: сервисы, а в них --- характеристики со своими правами}
\end{figure}

Пример ниже создаёт сервис с одной характеристикой «Яркость». Запишите в неё число 0--100 из
приложения nRF Connect на телефоне --- светодиод на \pin{4} изменит яркость.

\codecap{BLE-сервер с записываемой характеристикой яркости}
\codeinput{l12_ble}


\subsection{Разбор ошибки: сеть «подвисает»}

Не вызывайте сетевой код из ISR и не держите мьютекс во время Wi-Fi/MQTT-запроса. Обработчик кнопки
кладёт событие в очередь, а сетевая задача публикует его позднее. Для MQTT храните последнее состояние
и повторяйте подключение с задержкой. Для BLE проверяйте длину входного значения до копирования в буфер.
Так один неверный пакет не заблокирует всё устройство.

\begin{summary}
\begin{itemize}[leftmargin=*]
  \item ESP-NOW --- прямая короткая связь ESP; MQTT --- обмен через брокер; BLE --- близкая настройка.
  \item Радиоканал у ESP-NOW должен совпадать, а MAC-адреса и входные данные нужно проверять.
  \item Сетевые операции выполняются задачей с тайм-аутами, а не в ISR или бесконечном ожидании.
\end{itemize}
\end{summary}

\begin{selfcheck}
\begin{enumerate}
  \item Какой протокол выбрать для пульта без роутера, для умного дома и для настройки с телефона?
  \item Что означает retained-сообщение в MQTT?
  \item Почему сетевой код нельзя вызывать из обработчика прерывания?
\end{enumerate}
\end{selfcheck}

\begin{tasks}
\begin{enumerate}
\item \lv{1} Запустите ESP-NOW между двумя платами и убедитесь, что сообщение доходит.
\item \lv{1} Найдите MAC-адрес приёмника и объясните, зачем он нужен передатчику.
\item \lv{2} Сделайте ESP-NOW-пульт: короткое нажатие отправляет переключение, длинное --- яркость 100\%.
\item \lv{2} Опишите четыре MQTT-темы устройства: state, set, availability и diagnostics, и решите, какие из них retained.
\item \lv{3} Подключитесь к MQTT-брокеру: публикуйте значение потенциометра раз в 5 секунд и принимайте команду включения светодиода.
\item \lv{4} Спроектируйте BLE-сервис с характеристиками имени, яркости и версии прошивки; назначьте UUID и права доступа.
\end{enumerate}
\end{tasks}
